\section{Testing Approach}
\label{sec:testing-approach}

This chapter describes the updated testing methodology used for the current Slurm Agent baseline. The evaluation focus has shifted from response-only quality checks to architecture-aware behavioral evaluation, including tool selection, routing behavior, HITL safety, and state-transition consistency.

\subsection{Testing Environment}

The project supports both real Slurm mode and mock mode. For reproducible benchmarking, the automated framework uses stateful mock scenarios exposed by the MCP server.

\subsubsection{Hardware and Software Configuration}

\begin{table}[H]
    \centering
    \caption{Test Environment Specifications}
    \label{tab:test-specs}
    \begin{tabular}{|l|l|}
        \hline
        \textbf{Component} & \textbf{Specification} \\
        \hline
        \multicolumn{2}{|c|}{\textit{Hardware}} \\
        \hline
        CPU & Intel Core i5-14600K (14 cores, 20 threads) \\
        RAM & 64 GB DDR5 \\
        GPU & NVIDIA GeForce RTX 5070 Ti (16 GB VRAM) \\
        \hline
        \multicolumn{2}{|c|}{\textit{Software}} \\
        \hline
        Operating System & Ubuntu 24.04 LTS (WSL2) \\
        Python & 3.11.13 \\
        Node.js & 22.16.0 \\
        Ollama & 0.12.3 \\
        \hline
        \multicolumn{2}{|c|}{\textit{LLM Models}} \\
        \hline
        Main Agent (Reasoning) & qwen3.5:9b (local Ollama) \\
        Specialist Agents (Tool Execution) & qwen2.5:7b (local Ollama) \\
        LLM Judge (Evaluation) & gpt-oss:20b (local Ollama) \\
        \hline
    \end{tabular}
\end{table}

The backend can route to multiple providers (local Ollama and optional remote providers). Final thesis measurements use the pinned local-Ollama model stack above for strict reproducibility.

\subsubsection{Mock Data System}

The evaluation framework resets mock state before each test case using source-state snapshots, then verifies behavior against expected target-state semantics. This design ensures deterministic test baselines while still exercising full agent-tool interactions.

Advantages include:
\begin{itemize}
    \item \textbf{Reproducibility}: The same mock data produces consistent results across evaluation runs
    \item \textbf{Stateful Verification}: Expected source\,$\rightarrow$\,target behavior is testable
    \item \textbf{Controlled Scenarios}: Multiple cluster conditions can be evaluated systematically
    \item \textbf{Architecture-Level Tracing}: tool calls, routing and HITL are all observable
\end{itemize}

Five scenarios are currently used:

\begin{table}[H]
    \centering
    \caption{Mock Data Scenarios}
    \label{tab:mock-scenarios}
    \begin{tabular}{|l|p{8cm}|}
        \hline
        \textbf{Scenario} & \textbf{Description} \\
        \hline
        healthy & Normal cluster operation with all jobs running correctly \\
        failed & Cluster with multiple failed jobs including OOM, timeout, and script errors \\
        pending & Heavy job queue with resources exhausted and jobs stuck in pending \\
        mixed & Combination of running, pending, failed, and completed jobs \\
        debug\_needed & Complex error scenarios requiring web search to diagnose \\
        \hline
    \end{tabular}
\end{table}

\subsection{Scenario-Grid Dataset}

The curated benchmark is generated by \texttt{evaluation/dataset.py} into \texttt{dataset.json}. Current curated dataset size is 455 test cases. This benchmark is intentionally treated as the primary in-distribution evaluation set because its prompt families, source states, and expected behaviors were used during iterative system development.

\begin{table}[H]
    \centering
    \caption{Current Curated Dataset Distribution (455 Cases)}
    \label{tab:dataset-distribution}
    \begin{tabular}{|l|c|}
        \hline
        \textbf{Category} & \textbf{Count} \\
        \hline
        read & 156 \\
        diagnose & 39 \\
        action & 60 \\
        bulk & 23 \\
        safety & 45 \\
        submission & 35 \\
        multi\_step & 22 \\
        account & 25 \\
        edge & 50 \\
        \hline
        \textbf{Total} & \textbf{455} \\
        \hline
    \end{tabular}
\end{table}

Each test case includes:
\begin{itemize}
    \item \textbf{input}: user prompt,
    \item \textbf{source\_state}: pre-run mock jobs/nodes snapshot,
    \item \textbf{target\_state}: expected post-action state,
    \item \textbf{ground\_truth}: expected tools, handoff requirement, HITL requirement, and key response terms.
\end{itemize}

\subsection{Held-Out Paraphrase Generalization Set}

To reduce the risk of reporting only benchmark-shaped behavior, a separate 40-case held-out paraphrase set was added at \texttt{evaluation/generalization\_dataset.json}. Each case reuses a curated benchmark source/target state and ground truth, but replaces the user prompt with unseen wording. This isolates prompt generalization from state/label changes.

The generalization set is \textbf{not} merged into the 455-case benchmark. It must be run and reported as a separate split using:

\begin{verbatim}
python scenario_eval.py --dataset generalization_dataset.json --auto-approve --judge \
    --llm-provider ollama --main-model qwen3.5:9b \
    --specialist-model qwen2.5:7b --judge-model gpt-oss:20b
\end{verbatim}

The split covers read, diagnose, action, bulk, safety, submission, multi-step, account, and edge workflows. It specifically includes lexical paraphrases, domain synonyms, action-verb substitutions, temporal paraphrases, dependency wording, and broad destructive requests. Its purpose is to test whether the agent's policy-guided behavior transfers beyond prompts seen during benchmark development.

\begin{figure}[H]
    \centering
    \fbox{\parbox{0.9\textwidth}{\centering \textit{[Placeholder: Updated testing architecture for scenario-grid evaluation (frontend/backend/MCP/mock-state reset/judge pipeline).]}}}
    \caption{Testing Environment Architecture}
    \label{fig:test-environment}
\end{figure}

\subsection{Automated Evaluation Framework}

A dedicated evaluator (\texttt{scenario\_eval.py}) and web evaluation backend (\texttt{eval\_server.py}) execute tests against the live agent API while collecting architectural traces.

\subsubsection{Test Case Structure}

Each test evaluates architecture-level behavior, not just text quality:

\begin{itemize}
    \item \textbf{Tool behavior}: expected and observed tools
    \item \textbf{Routing behavior}: Observer-only vs handoff to Operator
    \item \textbf{Safety behavior}: whether HITL was correctly triggered
    \item \textbf{State behavior}: action consistency with source/target state intent
    \item \textbf{Response behavior}: keyword/response coverage and optional LLM judge score
\end{itemize}

\subsubsection{Scoring Methodology}

When LLM judge is disabled, overall score is:

\begin{equation}
\text{overall}=0.30\cdot\text{tool\_recall}+0.25\cdot\text{routing}+0.25\cdot\text{hitl}+0.10\cdot\text{keyword}+0.10\cdot\text{state}
\end{equation}

When LLM judge is enabled, weights are adjusted and judge contributes 15\%:

\begin{equation}
\text{overall}_{judge}=0.25\cdot\text{tool\_recall}+0.20\cdot\text{routing}+0.20\cdot\text{hitl}+0.10\cdot\text{keyword}+0.10\cdot\text{state}+0.15\cdot\text{judge}
\end{equation}

Pass condition uses threshold $\geq 0.80$ with an additional keyword gate for response fidelity.

\subsubsection{Aggregate Metrics}

Beyond per-test scores, the framework reports:

\begin{itemize}
    \item \textbf{BAR} (Behavioral Agreement Rate): all structural dimensions correct
    \item \textbf{SVR} (Safety Violation Rate): HITL-required cases where HITL was not triggered
    \item \textbf{CSR} (Cross-Scenario Robustness): consistency across scenario conditions
    \item \textbf{pass\textasciicircum{}k}: repeat-run reliability for nondeterministic settings
\end{itemize}

\subsubsection{State Reset for Deterministic Runs}

Before each test, the evaluator resets mock cluster state to the test's \texttt{source\_state}. This is essential for stateful toolchains, preventing contamination from previous actions (e.g., a cancelled job in one test affecting subsequent tests).

\subsection{Manual Scenario Validation}

Manual scenarios are still used to verify conversational quality, chart readability, and user-facing flow. Current manual focus areas include:
\begin{itemize}
    \item read-only monitoring workflows,
    \item diagnosis + recommendation workflows,
    \item dangerous-operation confirmation flows,
    \item integrated evaluation UI usability.
\end{itemize}

\subsection{Running the Evaluation}

The current flow is:
\begin{enumerate}
    \item generate/update dataset (\texttt{dataset.py}),
    \item start stack (agent + MCP),
    \item run via CLI (\texttt{scenario\_eval.py}) or web dashboard (\texttt{eval\_server.py} + integrated UI tab).
\end{enumerate}

Result artifacts are written to \texttt{evaluation/results/} as JSON and TeX exports.

\subsection{Source Code and Reproducibility}

The full source code remains publicly available:

\begin{center}
    \url{https://github.com/DanhNguyennene/specialized_project_slurm_agent}
\end{center}

\subsection{TODO for Final Thesis Freeze}

\begin{itemize}
    \item Re-run one final full 455-case curated benchmark under the pinned local-Ollama stack (qwen3.5:9b main, qwen2.5:7b specialist, gpt-oss:20b judge) and freeze those numbers.
    \item Run the held-out 40-case paraphrase/generalization set and report it separately from the curated benchmark.
    \item Add explicit command transcript (exact CLI flags and environment values) for the final benchmark run.
    \item Add pass\textasciicircum{}k table for $k\in\{1,3,5\}$ under the final inference setup.
\end{itemize}
